\parttitle{IX}{A Staged Repository: From NumPy Reservoir to Product}

This part designs, from scratch, a single repository that carries the whole
programme of this document: a small exact quantum reservoir in NumPy; the same
reservoir expressed as hardware-faithful circuits on an exact simulator; noisy
simulation of a simple Hamiltonian; a battery of more complex Hamiltonians;
qubit-reuse compilation of those Hamiltonians; an independent recurrence-free
reservoir (RF-QRC, Part~X); the integration of RF-QRC with the existing test
architecture and the reuse compiler; and finally the product pipeline of
Part~XI. The repository is not hypothetical: Stages~1--5 already exist as
verified code (the \code{code/} directory shipped with this document and the
companion qubit-reuse package), and the design below states exactly which
files move where, so that assembling the repository is a refactor with green
tests at every step rather than a rewrite.

\section{Design rules the repository inherits}
\label{sec:repo-rules}

Six rules, all earned earlier in this document or in the companion projects,
are promoted from habits to repository law.

\emph{R1 --- Anchors first.} No stage accepts code before its validation
anchors exist and fail meaningfully when sabotaged. The anchor suite of
Part~IV (inject/rebuild identity, propagator unitarity, $\rho$
Hermitian/trace-one/PSD, ridge recovery of a known linear map to $10^{-10}$,
$\NMSE(\text{mean})=1.000$, shot-noise $\sigma \propto S^{-1/2}$) seeds
\code{stage0\_anchors/}, and every later stage appends its own.

\emph{R2 --- The exact reference defines the ceiling.} The NumPy
density-matrix reservoir of Stage~1 is the semantic definition of every
architecture in the repository; framework ports (Qiskit, Cirq, PennyLane) are
validated \emph{against it at identical seeds}, never against intuition.

\emph{R3 --- One endianness point.} Qiskit count keys are little-endian;
\code{counts\_to\_features} performs the single \code{key[::-1]} reversal and
nothing else in the repository may reorder bits.

\emph{R4 --- Printed configuration, fixed seed.} Every experiment script
prints its full configuration block at runtime; \code{SEED = 7} project-wide;
no result enters a table unless a script printed it.

\emph{R5 --- Negative controls are first-class.} The all-to-all
zero-compression rows (SK, SYK$_4$), the Haar-random structureless reservoir,
and the mean-predictor $\NMSE=1.0$ line ship in the same harness as the
headline configurations and appear in the same tables.

\emph{R6 --- Promotion gates, not vibes.} A stage is \emph{done} when its
exit criterion --- a script with an exit code --- passes; downstream stages
import only from stages that are done. The repository's continuous
integration is exactly \code{run\_tests.py} run stage by stage.

\section{The stage ladder}
\label{sec:repo-ladder}

Each stage answers one question and owns one directory.
Table~\ref{tab:repo-stages} summarises; the subsections that follow give the
contents, the entry and exit criteria, and the mapping to existing assets.

\begin{table*}[h]
\centering\small
\caption{The repository's stage ladder. ``Exists'' names verified code that
seeds the stage; new files are set in italics.}
\label{tab:repo-stages}
\begin{tabular}{@{}clp{4.1cm}p{4.6cm}@{}}
\toprule
\# & Question answered & Seed assets & Exit criterion \\
\midrule
0 & do our instruments work? & anchor suites of both codebases & all anchors pass; sabotage test fails \\
1 & exact small-scale QRC (NumPy) & \code{qrc\_core.py}, \code{experiments.py} & NARMA-10/MC results reproduce Part~VII numbers \\
2 & same reservoir as circuits (exact) & \code{qiskit\_qrc\_hw.py}, \code{cirq\_qrc\_minimal.py} & circuit features $=$ NumPy features at seed 7 \\
3 & noisy simple Hamiltonian & Part~IV noise pipeline & noisy-vs-exact degradation curves printed; gate-local caveat stated \\
4 & complex Hamiltonians & seven-family benchmark (reuse repo) & $\langle r\rangle$ diagnostics hit 0.386/0.531 anchors; family scan runs \\
5 & qubit reuse on those & \code{qreuse\_*} (7 modules) & self-calibrated TVD gate; compression table incl.\ zero-compression rows \\
6 & independent RF-QRC & \emph{rfqrc\_reservoir/metrics/ baselines.py}, \emph{mfe\_model.py} & Part~X Phase~0--4 acceptance \\
7 & RF-QRC $\times$ reuse integration & \emph{qreuse\_batch.py} + Stage 5+6 & batch-and-compress validated; width independent of memory horizon \\
8 & product pipeline & \emph{data\_plane/ serving/ decision/ surfaces/ governance/} & Part~XI Phase~8--12 acceptance \\
\bottomrule
\end{tabular}
\end{table*}

\begin{strip}
\begin{lstlisting}[language={},caption={Repository tree. Stages 1--5 are populated by moving verified files; Stages 6--8 are populated by the plans of Parts X--XI.},label={lst:repo-tree}]
qrc-stack/
  README.md                 pyproject.toml            configs/
  run_tests.py              # anchors-first driver; exit 0 iff all stages green
  stage0_anchors/
    test_core_anchors.py    test_reuse_anchors.py     test_rfqrc_anchors.py
  stage1_numpy_core/
    qrc_core.py             baselines.py              tasks.py
    datasets.py             eda.py                    figstyle.py
  stage2_circuits_exact/
    qiskit_qrc_hw.py        cirq_qrc_minimal.py       pennylane_qrc_hybrid.py
  stage3_noisy_simple/
    noise_models.py         exp_tfim1d_noisy.py
  stage4_hamiltonian_battery/
    hamiltonians.py         exp_family_scan.py        # 7 families, <r> dial
  stage5_qubit_reuse/
    qreuse_ir.py            qreuse_analysis.py        qreuse_scheduler.py
    qreuse_compiler.py      qreuse_validation.py      qrc_experiment.py
    exp_reuse_compression.py
  stage6_rfqrc/
    rfqrc_reservoir.py      rfqrc_baselines.py        rfqrc_metrics.py
    mfe_model.py            exp_mfe_extremes.py       exp_lorenz63_check.py
  stage7_integration/
    qreuse_batch.py         exp_rfqrc_reuse.py        exp_vn_fusion_proto.py
    solar_ramps_data.py     exp_solar_ramps.py
  stage8_product/
    data_plane/   connectors.py qc.py transforms.py feature_store.py replay.py
    serving/      feature_provider.py heads.py conformal.py shadow_battery.py
    decision/     stacker.py regime_router.py thresholds.py alerts.py
    surfaces/     api.py dashboard/ agent/            # tools, guardrails, RAG
    governance/   monitors.py promotion_gates.py model_cards.py
  docs/                     # this handbook, regenerated figures
\end{lstlisting}
\end{strip}

\subsection{Stages 0--5: consolidating what exists}
\label{sec:repo-existing}

\emph{Stage 0.} The two anchor suites merge under one driver. The merged
driver keeps the qubit-reuse suite's self-calibration idea --- an equivalence
test passes only if TVD(direct, reuse) $\le 2.5\times$ the shot-noise floor
TVD(direct, direct-rerun) $+\,0.02$ --- and Part~IV's exact-simulation
anchors. A deliberate-sabotage mode (flip the endianness reversal, break a
Kraus normalisation) must turn the suite red; an anchor that cannot fail is
decoration.

\emph{Stage 1.} \code{qrc\_core.py} moves in unchanged: the Fujii--Nakajima
density-matrix reservoir \citep{fujii2017harnessing}, the classical baselines,
and the task generators. Its NARMA-10, memory-capacity, and ENSO numbers
(Part~VII) are the stage's regression targets. The recurrent 5-qubit CZ-ring
ENSO reservoir stays here as the long-memory comparator that Stage~6 must
face honestly.

\emph{Stage 2.} The framework ports (Part~IV) move in with their
validated-against-reference-at-identical-seed tests. This stage owns the
transpilation path (native basis, coupling maps) that makes ``hardware
simulator'' more than a phrase: circuits leave this stage ready for a device
even though this repository runs none.

\emph{Stage 3.} Noisy Aer on the simplest Hamiltonian family (TFIM-1D):
depolarising and thermal-relaxation models, degradation curves versus noise
strength and shots. The stage's documentation carries the caveat this project
measured the hard way: \emph{gate-local noise models cannot see the
$3$--$5\times$ depth inflation that qubit reuse introduces}, so Stage~5
noise claims must never be extrapolated from Stage~3 models alone.

\emph{Stage 4.} The seven-family Hamiltonian battery (TFIM-1D, mixed-field
Ising, XXZ, XY, TFIM-2D, Sherrington--Kirkpatrick, JW-mapped SYK$_4$) with
the disorder dial $W$ and level-spacing diagnostics
\citep{martinez2021dynamical,kobayashi2026edge}; the $\langle r\rangle$
anchors $0.386$ (Poisson) and $0.531$ (GOE) gate the stage.

\emph{Stage 5.} The reuse compiler's seven modules move in unchanged with
their compression scans. The stage's headline artefacts are the
interaction-graph law (nearest-neighbour chains compress $10\to5$ regardless
of gate count; the 2D lattice less; all-to-all not at all) and the
equivalence contract: reuse preserves the \emph{classical measurement
distribution} only --- exactly the contract a QRC read-out needs, and not one
a protocol needing the unmeasured state may assume (DeCross \emph{et al.},
Phys.\ Rev.\ X \textbf{13}, 041057 (2023)).

\subsection{Stages 6--8: the new work}
\label{sec:repo-new}

\emph{Stage 6} implements Part~X Phases~0--4: the recurrence-free reservoir
with virtual-node readout, its metric battery (predictability horizon,
event F-score versus prediction time, VPT, effective rank), the expanded
baseline set (NVAR, QELM ablation, Haar control), and the MFE extreme-event
anchor study, with the Lorenz-63 side check where the classical reservoir is
expected to win at matched degrees of freedom
\citep{ahmed2024prediction,ahmed2025optimal}.

\emph{Stage 7} implements Part~X Phases~5--6: the solar-ramp study on
registry data, the batch-and-compress path through the Stage-5 compiler
(\code{qreuse\_batch.py} lays $B$ independent step circuits side by side with
terminal measurements only, which the existing parser accepts today), the
structural claim that compiled width is independent of memory horizon
(contrast the windowed protocol's measured physical $=T{+}1$), and the
virtual-node fusion prototype behind an IR extension for
measurement-terminated blocks.

\emph{Stage 8} implements Part~XI Phases~8--12: the data plane with its
leakage firewall, the serving core behind the \code{FeatureProvider}
interface, the calibrated decision layer, the product surfaces including the
provenance-guarded agent, and the governance monitors whose promotion gates
are the pre-registered criteria of Parts~X--XI.

\section{Promotion protocol}
\label{sec:repo-promotion}

Work proceeds strictly up the ladder. To promote stage $k$: (i) its anchors
and exit script pass with exit code 0; (ii) its numbers reproduce the
regression targets of the documents that earned them (this handbook for
Stages~1--3, the reuse report for Stages~4--5, Parts~X--XI for Stages~6--8);
(iii) any interface it exports downstream is frozen and versioned; (iv) a
short stage report --- configuration block, tables, honest failures ---
lands in \code{docs/}. Cross-stage experiments (Stage~7 onwards) may only
import frozen interfaces. The repository is \emph{finished} in the same sense
this document is: a stranger rebuilds every figure and table with two
commands, and every claim is scoped theory, simulation, or hardware, with
surrogate data labelled wherever it appears.
